\subsection{The actual theta frame on a connected analytic family}
\label{subsec:ftb-actual-family-theta-frame}

We keep the supplied fine moduli object and the original tuple of
Proposition~\ref{prop:ftb-actual-local-pel-period-lift}.
In particular, \(E,H,D,l,X_o,\sigma,g,\mathrm{ty},D',t\) are
unchanged, \(1\leq D.p+D.q\), \(\sigma(\rho)=H.\mathrm{rt}_0\),
and \(\mathrm{ty}\) has the prescribed signature.
Let \(B\subset\mathbb C^m\) be open and preconnected.  On \(B\)
we have the actual algebraic points \(y_b\), the additive
surjective exponentials \(e_b:\mathbb C^g\to G_{y_b}(\mathbb C)\),
and the periods \(P_b:D'.\mathrm{La}\to\mathbb C^g\).
Keep all the conclusions of that proposition: their full
period kernels, both tangent-bijection clauses, the fixed
diagonal \(O_E\)-action, all positive-integer Weil identities
with the negative sign, both level markings, and joint
regular-evaluation holomorphy of \(y_b\) and \(e_b(z)\).
The periods \(P_b(a)\) are holomorphic for every \(a\).
This subsection uses the five algebraic laws retained in
Proposition~\ref{prop:ftb-full-fixed-fibre-theta-realization},
and the following precise additional algebraic
section/base-change law from \(\mathrm{Out.Laws}\).

For a principally polarized abelian scheme \(A/S\) with its
actual rigidified symmetric principal line \(L\), and \(S\)
locally Noetherian, each point of \(S\) has an open neighborhood
\(V\) on which the sections of \(L\) have a basis of size
\(1^g=1\), compatible with every affine base change
\(\operatorname{Spec}R\to V\).
Here the compatibility asserts that the pulled-back section
is an \(R\)-basis of the actual module of sections of the
pulled-back line, including for nonreduced \(R\).
We use this rank-one instance of
\(\mathrm{Out.SecBaseChangeSt}\), not merely the dimension of
the complex fibres.

Choose the source's integral basis
\[
 \mathrm{bas}:\mathbb Z^g\oplus\mathbb Z^g
                  \xrightarrow{\ \sim\ }D'.\mathrm{La},
 \qquad
 D'.EF(\mathrm{bas}(v),\mathrm{bas}(w))
                  =s\,\omega_{\rm std}(v,w),\quad s\in\{1,-1\}.
 \tag{TF1}
\]
Put \(p_i(b)=P_b(\mathrm{bas}(e_i,0))\) and
\(q_i(b)=P_b(\mathrm{bas}(0,e_i))\).
Suppose matrices \(h_b,\tau_b\) satisfy
\[
 h_bp_i(b)=e_i,\qquad h_bq_i(b)=\tau_be_i,\qquad
                    \operatorname{Im}\tau_b>0 .
 \tag{TF2}
\]
The function \(\theta_{\alpha,\beta}(z,\tau)\) below is the
source's classical theta function with characteristic;
its two characteristic vectors lie in \(\{0,\tfrac12\}^g\).

\begin{lemma}[The half-characteristic on the actual fixed fibre]
\label{lem:ftb-actual-half-characteristic}
For each \(b\in B\), there is a unique pair
\((\alpha_b,\beta_b)\in\{0,\tfrac12\}^g\times
\{0,\tfrac12\}^g\) with the following property.
The actual line \(L_{y_b}\) admits nonvanishing holomorphic
frame coordinates \(f_{b,j}(z)\), on the inverse images under
\(e_b\) of its original algebraic charts, such that
\[
 \begin{split}
 f_{b,j}(z)&=g_{jj'}(e_b(z))f_{b,j'}(z),\\
 f_{b,j}(0)&=\text{the value of the given rigidification
                       in chart \(j\)},\\
 s_{b,j}(e_b(z))&=
       \theta_{\alpha_b,\beta_b}(h_bz,\tau_b)f_{b,j}(z)
 \end{split}
 \tag{TF3}\label{eq:ftb-actual-family-theta-frame}
\]
for one actual algebraic section \(s_b\in H^0(G_{y_b},L_{y_b})\).
Uniqueness refers to the pair of half-characteristics; it does
not assert uniqueness of the analytic frame.
\end{lemma}
\begin{proof}
Pull the actual group, cubical line, rigidification and symmetry
to the fibre \(y_b\).  Properness, relative ampleness and
triviality of the polarization kernel pull back, so this is
the principally polarized abelian variety to which
Proposition~\ref{prop:ftb-full-fixed-fibre-theta-realization}
applies.  The pullback coordinates and their transition units
are the original ones: evaluating a pulled-back transition
at the lifted point gives \(g_{jj'}(e_b(z))\).
Thus a fibre frame can subsequently be transported back to
the coordinates in (TF3), including the origin
rigidification.

Let \(u=(e_u,\Lambda_u,f_{u,j})\) be the resulting good
uniformization, with its positive Hermitian form \(H_u\),
semicharacter \(\chi\), and section correspondence.
The source's linear-lift argument gives a complex linear
isomorphism \(F\) with \(e_uF=e_b\).
Indeed \(e_b\) has countable kernel \(P_b(D'.\mathrm{La})\).
A kernel of the linear lift is a vector subspace contained
in that countable set and is therefore zero.
Surjectivity modulo the countable lattice \(\Lambda_u\)
forces its real cokernel to be zero as well; this is exactly
the countable-kernel/countable-cokernel step in
Lemma~\ref{lem:ftb-holomorphic-linear-lift}.
Consequently
\[
 Fz\in\Lambda_u\ \Longleftrightarrow\
                  z\in P_b(D'.\mathrm{La}).
\]
Compare the actual Weil relations for \(e_u\) and \(e_b\)
through this isomorphism of points and the pulled-back
line.  For every \(k>0\), their values are respectively
\(\exp(-2\pi i\,\operatorname{Im}H_u(FP_ba,FP_ba')/k)\)
and \(\exp(-2\pi i\,D'.EF(a,a')/k)\).
The two exponents are integers.  They agree for every \(k\),
so the integers agree: take \(k\) larger than the absolute
value of their difference.  Hence
\[
 \operatorname{Im}H_u(FP_ba,FP_ba')=D'.EF(a,a').       \tag{TF4}
\]

The first identity in (TF2) makes \(h_b\) invertible.
In these coordinates
\[
 h_bP_b(\mathrm{bas}(m,n))=m+\tau_bn\quad(m,n\in\mathbb Z^g).
\]
Pull \(H_u\) through \(Fh_b^{-1}\).
The first, mixed and last blocks of (TF4), using (TF1),
are the elementary Riemann-form relations of
the source's normal-form calculation.
They show that \(\tau_b\) is symmetric and that this
positive Hermitian matrix is
\((\operatorname{Im}\tau_b)^{-1}\).
For clarity about the sign, the real matrix \(Y\) of that
form satisfies \(Y\operatorname{Im}\tau_b=-sI\).
When \(g>0\), positivity of both matrices rules out \(s=1\):
the same normal-form relations would make
\(\operatorname{Im}(-\tau_b)\) positive definite.
Thus \(s=-1\).  In dimension zero all these identities
are vacuous.  This is the source's
\(\mathrm{isSymm\_and\_eq\_imInv}\) calculation.

We next obtain signs for the semicharacter of the actual
symmetric line.  Write its Appell--Humbert multiplier as
\[
 J_\lambda(z)=\chi(\lambda)
    \exp\{\pi H_u(z,\lambda)+\pi H_u(\lambda,\lambda)/2\},
 \qquad f_{u,j}(z)=J_\lambda(z)f_{u,j}(z+\lambda).
\]
The laws for \(\lambda,-\lambda\), applied successively,
give \(\chi(\lambda)\chi(-\lambda)=1\).
Evaluate the actual symmetry isomorphism
\([-1]^*L\simeq L\) in two line charts, obtaining
holomorphic units \(T_{ij}(z)\).
The transition law makes
\[
             R(z)=T_{ij}(z)\frac{f_{u,j}(z)}{f_{u,i}(-z)}
\]
independent of the selected pair of charts.
It is an entire unit, and its multiplier under \(\lambda\)
is \(\chi(-\lambda)/\chi(\lambda)\).
The Hermitian part of this multiplier is zero.
Lemma~\ref{lem:ftb-trivial-appell-humbert-data}, applied with
zero Hermitian form, therefore makes this character trivial.
It follows that \(\chi(-\lambda)=\chi(\lambda)\), and
\(\chi(\lambda)^2=1\).

Choose \(\alpha_i,\beta_i\in\{0,\tfrac12\}\) whose exponential
signs are the values of \(\chi\) on the first and second
integral period generators.  The source verifies the full
multiplier law from these \(2g\) signs as follows.
For the desired theta function let \(S\) be the set of lattice
elements on which that multiplier law holds.
Nonvanishing of \(J_\lambda\) and its cocycle identity make
\(S\) an additive subgroup: multiplication proves closure
under addition, and the laws for \(\lambda,-\lambda\)
prove closure under negation.
The explicit theta quasi-period formulas put every basis
generator in \(S\).  Those generators span the lattice, so
the law holds for all its elements.
The Gaussian-normalized function \(\mathrm{canTh}\) is
therefore an analytic theta function for \((H_u,\chi)\).
The already proved section correspondence supplies its
actual algebraic section.

Set
\[
 Q(\zeta)=\exp\{\tfrac\pi2\,
                  \zeta^t(\operatorname{Im}\tau_b)^{-1}\zeta\}.
\]
In the source's convention
\(\mathrm{canTh}(\zeta)=Q(\zeta)
\theta_{\alpha,\beta}(\zeta,\tau_b)\).
Thus the frames
\(Q(h_bz)f_{u,j}(Fz)\) give (TF3).
Their origin values are the original ones because \(Q(0)=1\).

Here are the uniqueness details, including the unit ambiguity
in an analytic frame.
Two nonvanishing frames of the same pulled-back line have a
nowhere-zero entire ratio \(r\), because their ratios in line
charts agree by cancellation of \(g_{jj'}\).
A second characteristic and such a frame yield another
analytic theta function for \((H_u,\chi)\).
The Fourier coefficient bound for the principal lattice in
degree one makes that analytic theta space one dimensional.
The first theta function is nonzero, so the two are scalar
multiples; the scalar is nonzero.
After undoing \(Q\) and the frame ratio this gives
\[
 \theta_{\alpha',\beta'}(\zeta,\tau_b)
                   =U(\zeta)\theta_{\alpha,\beta}(\zeta,\tau_b)
                                                        \tag{TF5}
\]
with \(U\) an entire unit.

The source reduces (TF5), using the theta shift formula, to
\[
 \theta_{0,0}(\zeta+w,\tau_b)
                    =V(\zeta)\theta_{0,0}(\zeta,\tau_b),
 \quad w=\tau_b(\alpha'-\alpha)+(\beta'-\beta),
\]
with another entire unit \(V\).
We recall its translate test.  The quasi-period formulas and
cancellation on the nonzero open set of theta, followed by
the identity theorem, give
\[
 V(\zeta+\tau_bn+m)=e^{-2\pi i\,n^tw}V(\zeta).
\]
Its logarithmic derivative \(V^{-1}dV\) is an entire
lattice-periodic linear-form-valued function, hence constant
by the fundamental-parallelepiped bound and Liouville.
Integrating the zero derivative of \(V e^{-\ell}\) gives
\(V(\zeta)=V(0)e^{\ell(\zeta)}\).
Integer periods imply
\(\ell(e_i)=2\pi i\,k_i\) with \(k_i\in\mathbb Z\);
the \(\tau_be_i\) periods then imply
\(w+\tau_bk\in\mathbb Z^g\).
Thus \(w\in\mathbb Z^g+\tau_b\mathbb Z^g\).
Taking imaginary parts in
\(\tau_b(\alpha'-\alpha)+(\beta'-\beta)=\tau_bn+m\)
and using invertibility of \(\operatorname{Im}\tau_b\)
gives \(\alpha'-\alpha=n\).
Each coordinate difference has absolute value at most
\(1/2<1\), so \(n=0,\alpha'=\alpha\), and then likewise
\(m=0,\beta'=\beta\).
This proves uniqueness without assuming that a frame ratio
was constant at the outset.
\end{proof}

\begin{lemma}[A local algebraic section that spans every actual fibre]
\label{lem:ftb-local-algebraic-theta-section}
For each parameter \(b_0\), there is an algebraic open
\(V\subset X_o\) containing \(y_{b_0}\) and a section
\(s^0=(s^0_j)\) of the actual \(L\) over \(G|_V\) which is
nonzero on every fibre over \(V\).
Every algebraic section of \(L_{y_b}\), for \(y_b\in V\),
is a scalar multiple of this restriction.
\end{lemma}
\begin{proof}
Apply the rank-one section/base-change law stated above.
On \(G|_V\) its one basis section is initially written in
the pullback charts.  Transport those coordinates through
the open immersion \(G|_V\to G\).
Its stalk maps are isomorphisms, hence injective, so the
transition identities descend to the original charts and
give \(s^0\).  Under any affine base change the same
coordinates are the given basis section: the equality
is checked on pulled-back chart germs, using the section
pullback identity.  In particular they span every complex
fibre.

The section cannot vanish on a fibre over a point \(v\in V\).
Otherwise all its chart coordinates on that residue-field
fibre have empty nonzero loci.  The fibre is smooth over
its residue field and is reduced.  Every such coordinate
would therefore be zero.  Applying the linear independence
clause of the rank-one basis after that base change would
give \(1=0\) in the residue field, a contradiction.
The resulting nonzero open in the fibre has a complex
point after a complex base change, by the finite-type
closed-point theorem.  Thus the restriction is genuinely
nonzero as a section on each of the actual complex fibres.
Equivalently, one can use the fibre dimension-one law:
the value line of all sections has dimension one, and
evaluation at a point where \(s^0\ne0\) makes its scalar
nonzero.  Both descriptions give the asserted spanning
in the original line coordinates.
\end{proof}

\begin{proposition}[One characteristic and jointly holomorphic frames]
\label{prop:ftb-actual-family-theta-frame}
There are fixed \(\alpha,\beta\in\{0,\tfrac12\}^g\), chosen
once on \(B\), for which (TF3) holds on every fibre, and
the frame coordinates can be chosen jointly holomorphic
in \((b,z)\) on their actual inverse-image line charts.
Their transition functions and their origin values are
the original transition functions and rigidification.
Thus this is the source's
\(\mathrm{IsThetaFrame}\) for the same universal line,
points and exponentials.
\end{proposition}
\begin{proof}
First the displayed theta functions vary holomorphically.
If \(M_b\) is the matrix with columns \(p_i(b)\), then
\(h_bM_b=I\).  The adjugate formula makes \(h_b=M_b^{-1}\)
holomorphic.  The second identity of (TF2) then makes
\(\tau_b\) holomorphic.  Local normally convergent theta
series on the positive-imaginary domain give joint
holomorphy of every fixed half-characteristic function
\[
             \Theta_c(b,z)=\theta_c(h_bz,\tau_b).
\]
It is not identically zero on any fibre: theta is nonzero
and \(h_b\) is onto.  The identity theorem therefore gives
frequent nonvanishing in every fibre neighborhood,
including at its zeros.

By Lemma~\ref{lem:ftb-actual-half-characteristic} each fibre
selects a unique characteristic \(c(b)\) from the finite
set \(\{0,\tfrac12\}^{2g}\).
We show that every choice locus is closed relative to \(B\).
Fix \(b_0\in B\cap\overline{\{b\in B:c(b)=c\}}\).
Choose \(s^0,V\) as in
Lemma~\ref{lem:ftb-local-algebraic-theta-section}.
On an actual line chart around \((b_0,z_0)\) its reading
\[
                 f_j(b,z)=s^0_j(e_b(z))
\]
is jointly holomorphic.
At parameters where \(c(b)=c\), fibre spanning gives a
nonzero scalar \(a_b\) with
\[
       f_j(b,z)=a_b^{-1}f_{b,j}(z)\Theta_c(b,z).       \tag{TF6}
\]
Consequently both \(f_j/\Theta_c\) and
\(\Theta_c/f_j\) extend holomorphically in that fibre
chart and are reciprocal units.

Apply the circle-quotient construction of
Lemma~\ref{lem:ftb-circle-quotient} in both directions.
For the first direction use frequent nonvanishing of
\(\Theta_c(b_0,-)\).
For the second use the selected frame at \(b_0\):
its relation
\(f_j=a_{b_0}^{-1}f_{b_0,j}\Theta_{c(b_0)}\)
and nonvanishing of \(f_{b_0,j}\) imply frequent
nonvanishing of \(f_j(b_0,-)\).
We obtain holomorphic local quotients \(G_1,G_2\) at
the limiting parameter, with
\[
 G_1\Theta_c=f_j,\qquad G_2f_j=\Theta_c .
\]
On the frequent nonzero set of \(\Theta_c\) these give
\(G_2G_1=1\); continuity gives the same identity at the
chosen zero as well.
Hence
\[
 \rho(z)=G_2(b_0,z)\,a_{b_0}^{-1}f_{b_0,j}(z)
\]
is a local holomorphic unit with
\(\rho\,\Theta_{c(b_0)}=\Theta_c\).
These local quotients agree on overlaps: cancel the
frequently nonzero \(\Theta_{c(b_0)}\) and then use
continuity.  They glue to an entire unit \(\rho\).
Rescale the selected frame at \(b_0\) by
\(\rho(0)/\rho(z)\) and its algebraic section by
\(\rho(0)\).  This gives a frame of the actual same
line for \(\Theta_c\), preserving the origin
rigidification.  Uniqueness in
Lemma~\ref{lem:ftb-actual-half-characteristic} now gives
\(c(b_0)=c\).

A map to a finite set with these relatively closed fibres
is locally constant: the complement of any fibre is a
finite union of the other closed fibres.
Preconnectedness makes it constant on \(B\).
If \(B\) is empty choose the zero characteristic and all
assertions are vacuous.

It remains to make the frames jointly holomorphic; an
arbitrary fibrewise choice does not by itself do this.
Choose fibre frames \(f_{b,j}\) for the now fixed
characteristic, and near \(b_0\) write (TF6) with that
characteristic and the local section \(s^0\).
Put \(F_j(b,z)=a_b^{-1}f_{b,j}(z)\).
Each \(F_j\) is holomorphic in the fibre variable and
satisfies \(F_j\Theta_c=f_j\).
The second assertion of
Lemma~\ref{lem:ftb-circle-quotient} makes it jointly
holomorphic.

Choose a line chart \(j_0\) containing the identity near
\(b_0\), and let \(r_{j_0}(b)\ne0\) be the holomorphic
value of the actual rigidification in that chart.
Since
\[
          F_{j_0}(b,0)=a_b^{-1}r_{j_0}(b)\ne0,
 \qquad
          a_b=r_{j_0}(b)/F_{j_0}(b,0),
\]
the scalar \(a_b\) is holomorphic near \(b_0\).
This argument does not divide by theta at the origin,
which may be zero.
Therefore \(f_{b,j}=a_bF_j\) is jointly holomorphic.
The transition, nonvanishing, rigidification and
actual algebraic fibre-section identities already hold
for these fibre frames; they give all clauses of
\(\mathrm{IsThetaFrame}\).
\end{proof}

\paragraph{Frozen-source correspondence and retained inputs.}
All paths below have prefix
\texttt{src/Development/AbelianVarietyTheory/Proofs/}.
The sign argument is
\texttt{CubicalLineBundles/SemicharacterSqEqOneOfSymmetric.lean}
(\texttt{d2d717aec9d26a9ed8ca968455911d6367e362bf}),
\srcid{semicharacter_sq_eq_one_of_symmetric}.
The basis-sign construction, multiplier subgroup and
characteristic uniqueness are
\texttt{ModuliProofs/ExistsHalfOfSqEqOne.lean}
(\texttt{23a80692419984a24014f7042cef8bee1ddfed8c}),
\srcid{exists_half_char}, \srcid{half_char_unique};
the translate-unit test is
\texttt{ToricGeometry/MemPeriodLatticeOfThetaTranslateEqUnitMul.lean}
(\texttt{eed4913226755b63075a8f4dd698a5aa11894c89}).
The actual fibre-frame and Gaussian transport are
\texttt{ModuliProofs/ExistsUniqueHalfChar.lean}
(\texttt{a5b282df7dafc2ba7da9c949db64ff79d4f6a38d}),
\srcid{exists_unique_half_char}, and
\texttt{ModuliProofs/ExistsUniqueThetaCharAtPoint.lean}
(\texttt{3bb3de1f849094a745940ee6c080668be700ceae}),
\srcid{existsUnique_char_of_triple}.
Its public wrapper
\texttt{ModuliProofs/ExistsUniqueThetaCharOfUnif.lean}
(\texttt{82e1a4972cb06562e49d26f54b88952d80c799fe})
uses that same producer.

The local section construction is
\texttt{UnitaryExistence/ExistsNonvanishing.lean}
(\texttt{b717c3b20851284c2e7162d10270f736f6bf981d}),
\srcid{PushInv.exists_basisOn},
\srcid{PushInv.exists_nonvanishing}, and
\srcid{PushInv.pushInvertible};
the rank-one fibre value-line comparison is
\texttt{ModuliProofs/ExistsValuesLineFibre.lean}
(\texttt{048b56a4292b1ed2c08c8dc28099fda1f3f31a2e});
and its local nonzero-section wrapper is
\texttt{ModuliProofs/ExistsSecOnNonvanishing.lean}
(\texttt{d3c5d67c00d7537db821abd5bd68fbc2958c09ac}).
The two reciprocal quotients, entire-unit gluing and
finite closed-fibre test are
\texttt{ModuliProofs/EqOfFiniteOfClosedFibres.lean}
(\texttt{7f7959eed9e909e48e267a075fa077ab74d989e9})
and
\texttt{ModuliProofs/ExistsFibreFrameOfMemClosure.lean}
(\texttt{b97e854606513593d3fa2ebbae053163373bd288}),
including \srcid{exists_fibreFrame_of_mem_closure} and
\srcid{isThetaFrame_of_fibreFrames}.
Their final assembly is
\texttt{ModuliProofs/ExistsThetaFrameOfPreconnected.lean}
(\texttt{6f0015a5b36bb1ce80e70e405023a80a95159307})
and \texttt{ModuliProofs/ExistsThetaFrameFamily.lean}
(\texttt{1e14b15190f65d3d76f3a98fa378b536161565fc}),
\srcid{exists_thetaFrame_family}.

The conclusion is the actual analytic-family theta-frame
producer, for the supplied fine moduli object and its
original universal line.
It does not by itself construct that moduli object,
prove the algebraic section/base-change law, algebraize
an arbitrary positive complex torus, or construct the
common finite Weil representation and its normalized
vacuum.  Those are distinct source suppliers.
